% !TEX root = ../main.tex
\section*{Executive Summary}
\addcontentsline{toc}{section}{Executive Summary}

\plantodo{S14, S16}{Write this section LAST. Distill: (1) the improvement
opportunity/problem statement, (2) the MBSE approach used, (3) major
architecture and optimization-demonstration findings stated quantitatively,
(4) the operational impact/benefit of the reference architecture (see
Section~\ref{sec:impact}), and (5) recommendations. Limit to
1-2 pages; no citations, undefined terms, or content not found elsewhere in
the report. State the two claims here as well as in the Introduction
(Section~\ref{sec:introduction}; owner's note on the 2026-10-11 print, page 11).}

\begin{itemize}
    \item \textbf{Problem.} NAS decisions are made by independent actors, each with partial information. When a conflict or disruption occurs, the decision is made quickly and myopically. The sector controller resolves conflicts knowing type, level, speed and route, but not the aircraft's performance.
    \item \textbf{Gap.} Few sources model the NAS with systems engineering. None found address what detailed aircraft state information is worth to the decision.
    \item \textbf{Claims.}
    \begin{itemize}
        \item Experiment: a higher-fidelity decision aid, fed accurate performance data and the wider airspace state, resolves conflicts better, as measured against the stakeholder objectives the study keeps.
        \item Model: an MBSE architecture locates the missing interface, names who owns the decision and the data, states the unmet requirement, and lets the scoring fidelity be swapped with the structure unchanged.
    \end{itemize}
    \item \textbf{Approach.}
    \begin{itemize}
        \item SysML v2 model of the NAS as a system of systems: context diagram, internal block diagrams, nominal-operations activity, requirements, verification cases, the proposed airline-to-ATC interface.
        \item Experiment: one level-crossing conflict, four aircraft, one sector; four cases (A to D) with increasing information; options scored with OpenAP at three fidelities; impacts reported side by side for aircraft, airline, sector and system.
    \end{itemize}
    \item \textbf{Findings: architecture.} [Measures from Section~\ref{sec:results}: requirements traced, requirements allocated, requirements tagged \texttt{improved}; the gap the model flags from the architecture alone (Section~\ref{sec:model-gap}); the verdicts by case (Appendix~\ref{app:thread-trace}).]
    \item \textbf{Findings: experiment.} [Option chosen by each case; B minus A (value of automation) and C minus B (value of the information exchange) in fuel, CO2 and instructions; who pays.]
    \item \textbf{Impact.} [From Section~\ref{sec:impact}: operational and dollar terms where a credible conversion exists.]
    \item \textbf{Recommendations.} [From the Conclusions and Recommendations section.]
\end{itemize}

Decisions made within the United States National Airspace System (NAS) are often made by independent actors, each with only partial information, without a comprehensive understanding of the decisions impact on the other systems and actors involved, nor the stakeholders that are affected. From instantaneous operational control decisions to long-term strategic planning, the lack of a holistic view can lead to suboptimal outcomes. This report identifies and investigates one significant gap in the current NAS decision-making process: the absence of detailed aircraft state information in conflict resolution scenarios. When a conflict or disruption occurs, the decision is made quickly and myopically. The sector controller resolves conflicts knowing type, level, speed and route, but not the aircraft's performance. A higher-fidelity decision aid, fed accurate performance data and the wider airspace state, resolves conflicts better, as measured against the stakeholder objectives.

The identification of the gap in the current NAS is not new, but the approach taken in this report is novel. Few sources model the NAS with systems engineering, and none found address what detailed aircraft state information is worth to NAS decisions. This report presents a model-based systems engineering (MBSE) architecture that locates the missing interface, names who owns the decision and the data, states the unmet requirement, and allows for scoring fidelity to be swapped with the structure unchanged. In doing so, the report demonstrates the value of MBSE in identifying and addressing gaps in complex systems like the NAS.

The reference architecture developed in this report is based on a SysML v2 model of the NAS as a system of systems. The model includes a context diagram, internal block diagrams, activity diagrams, requirements, verification cases, and the proposed airline-to-ATC interface. An analysis of the model reveals the gap in the current NAS decision-making process, highlighting the need for a higher-fidelity decision aid that incorporates detailed aircraft state information. The model also provides a framework for evaluating the impact of different levels of information on conflict resolution outcomes.

An experiment is conducted that involves one level-crossing conflict with four aircraft in one sector, with four cases (A to D) representing increasing levels of information. The options are scored using an open source aircraft performance model tool serving as the truth source. Three model fidelities are assessed and the impacts are reported side by side for aircraft, airline, sector, and system.

The experiment results demonstrate that the higher-fidelity decision aid, which incorporates detailed aircraft state information, leads to better conflict resolution outcomes. The findings show that the value of automation and the value of the information exchange can be quantified in terms of fuel savings, CO2 emissions reduction, and the number of instructions issued. The report also identifies who bears the costs and benefits of these improvements.

